% =========================================================================
\section{Sistemas de coordenadas y operadores diferenciales}\label{sec:coordenadas}
% =========================================================================
Las leyes del electromagnetismo no dependen de los ejes que elijamos, pero los
\emph{cálculos} sí: ante una distribución con simetría esférica o cilíndrica,
lo inteligente es usar coordenadas adaptadas a esa simetría. En esta sección
repasamos los tres sistemas ortogonales fundamentales (cartesiano, cilíndrico y
esférico) con sus elementos de línea, de área y de volumen, sus cambios de base,
y los cuatro operadores diferenciales en cada uno de ellos~\cite{FORMULARIO}.
Los obtenemos todos a la vez a partir de un marco general.

% -------------------------------------------------------------------------
\subsection{Coordenadas curvilíneas ortogonales}\label{sec:curvilineas}
% -------------------------------------------------------------------------
Un sistema de coordenadas curvilíneas $(q_1,q_2,q_3)$ viene dado por el
vector de posición $\vc{r}=\vc{r}(q_1,q_2,q_3)$. Variando solo $q_i$ se
recorre la \emph{línea coordenada} $i$, cuya tangente es $\partial\vc{r}/\partial q_i$.

\begin{definicion}[Coordenadas ortogonales y factores de escala]\label{def:ortogonales}
El sistema es \emph{ortogonal} si los tres vectores $\partial\vc{r}/\partial q_i$
son perpendiculares dos a dos. Se definen entonces los \emph{factores de escala}
$h_i$ y los versores $\uh{i}$ por
\[
  h_i\defeq\Bigl\lVert\pd{\vc{r}}{q_i}\Bigr\rVert,
  \qquad
  \uh{i}\defeq\frac{1}{h_i}\pd{\vc{r}}{q_i},
\]
y $(\uh{1},\uh{2},\uh{3})$ es una base ortonormal \emph{local} (depende del
punto), que supondremos orientada a derechas: $\uh{1}\times\uh{2}=\uh{3}$.
\end{definicion}

El factor $h_i$ convierte un incremento de la coordenada, $\dif q_i$, en una
\emph{longitud} $h_i\,\dif q_i$ (por ejemplo, $\dif\varphi$ es un ángulo y
$\rho\,\dif\varphi$ un arco). De ahí se obtienen los elementos geométricos.

\begin{proposicion}[Elementos de línea, área y volumen]\label{prop:elementos}
En coordenadas ortogonales,
\begin{align}
  \dl&=h_1\,\dif q_1\,\uh{1}+h_2\,\dif q_2\,\uh{2}+h_3\,\dif q_3\,\uh{3}, \label{eq:dl-gen}\\
  \dif S_i&=h_jh_k\,\dif q_j\,\dif q_k\quad(i,j,k\ \text{cíclicos}), \label{eq:dS-gen}\\
  \dV&=h_1h_2h_3\,\dif q_1\,\dif q_2\,\dif q_3, \label{eq:dV-gen}
\end{align}
donde $\dif S_i$ es el área de la cara del elemento de volumen perpendicular a
$\uh{i}$ (a $q_i$ constante). El jacobiano del cambio de variable
$(q_1,q_2,q_3)\mapsto(x,y,z)$ es
$\bigl|\partial(x,y,z)/\partial(q_1,q_2,q_3)\bigr|=h_1h_2h_3$.
\end{proposicion}
\begin{proof}
Al variar $q_i$ en $\dif q_i$ el punto se desplaza
$(\partial\vc{r}/\partial q_i)\dif q_i=h_i\,\dif q_i\,\uh{i}$, y el desplazamiento
total es la suma, \cref{eq:dl-gen}. El elemento de volumen es el paralelepípedo
rectangular de aristas $h_i\,\dif q_i$, lo que da \eqref{eq:dV-gen}, y cada cara es
el rectángulo de las otras dos aristas, \cref{eq:dS-gen}. El jacobiano es el
determinante de la matriz cuyas columnas son $\partial\vc{r}/\partial q_i=h_i\uh{i}$,
que, por ser ortogonales, vale el producto de sus longitudes.
\end{proof}

\begin{teorema}[Operadores diferenciales en coordenadas ortogonales]\label{teo:operadores-ortogonales}
Sean $f$ un campo escalar y $\vc{A}=A_1\uh{1}+A_2\uh{2}+A_3\uh{3}$ un campo vectorial.
Con $\partial_i\equiv\partial/\partial q_i$ y $H=h_1h_2h_3$,
\begin{align}
  \nab f&=\sum_{i=1}^{3}\frac{1}{h_i}\,\partial_if\;\uh{i}, \label{eq:grad-gen}\\
  \nab\cdot\vc{A}&=\frac{1}{H}\Bigl[\partial_1(h_2h_3A_1)+\partial_2(h_3h_1A_2)+\partial_3(h_1h_2A_3)\Bigr], \label{eq:div-gen}\\
  \nab\times\vc{A}&=\frac{1}{H}
  \begin{vmatrix}
    h_1\uh{1} & h_2\uh{2} & h_3\uh{3}\\
    \partial_1 & \partial_2 & \partial_3\\
    h_1A_1 & h_2A_2 & h_3A_3
  \end{vmatrix}, \label{eq:rot-gen}\\
  \Lap f&=\frac{1}{H}\Bigl[\partial_1\Bigl(\frac{h_2h_3}{h_1}\partial_1f\Bigr)
     +\partial_2\Bigl(\frac{h_3h_1}{h_2}\partial_2f\Bigr)
     +\partial_3\Bigl(\frac{h_1h_2}{h_3}\partial_3f\Bigr)\Bigr]. \label{eq:lap-gen}
\end{align}
\end{teorema}
\begin{proof}
\emph{Gradiente.} Por un lado $\dif f=\sum_i\partial_if\,\dif q_i$; por otro, por
la \cref{eq:df-compacta} y la \cref{eq:dl-gen},
$\dif f=\nab f\cdot\dl=\sum_i(\nab f)_i\,h_i\,\dif q_i$. Comparando coeficientes de los
$\dif q_i$ independientes, $(\nab f)_i=\partial_if/h_i$.

\emph{Divergencia.} Usamos la \cref{prop:div-intrinseca} sobre el elemento de
volumen curvilíneo. El flujo por las dos caras perpendiculares a $\uh{1}$ (de
área $h_2h_3\,\dif q_2\,\dif q_3$, evaluada en $q_1$ y en $q_1+\dif q_1$) es
\[
  \bigl[A_1h_2h_3\bigr]_{q_1+\dif q_1}\dif q_2\dif q_3-\bigl[A_1h_2h_3\bigr]_{q_1}\dif q_2\dif q_3
  \approx\partial_1(h_2h_3A_1)\,\dif q_1\dif q_2\dif q_3 ,
\]
y análogamente para las otras dos parejas de caras. Dividiendo el flujo total
por $\dV=H\,\dif q_1\dif q_2\dif q_3$ se obtiene \eqref{eq:div-gen}.

\emph{Rotacional.} Usamos la \cref{prop:vorticidad} sobre la cara perpendicular a
$\uh{1}$, de área $h_2h_3\,\dif q_2\,\dif q_3$. Su circulación (lados de longitud
$h_2\dif q_2$ y $h_3\dif q_3$) es, a primer orden,
$\bigl[\partial_2(h_3A_3)-\partial_3(h_2A_2)\bigr]\dif q_2\dif q_3$; dividiendo por el
área, $(\nab\times\vc{A})_1=\tfrac{1}{h_2h_3}\bigl[\partial_2(h_3A_3)-\partial_3(h_2A_2)\bigr]$.
Las otras dos componentes se obtienen permutando cíclicamente los índices, y el
conjunto es el desarrollo del determinante \eqref{eq:rot-gen}.

\emph{Laplaciano.} Es $\nab\cdot(\nab f)$: basta sustituir $A_i=\partial_if/h_i$
en la \cref{eq:div-gen}.
\end{proof}

Los factores de escala de los tres sistemas que usaremos son
\begin{center}
\renewcommand{\arraystretch}{1.3}
\begin{tabular}{@{}lccc@{}}
\toprule
Sistema & $(q_1,q_2,q_3)$ & $(h_1,h_2,h_3)$ & $H=h_1h_2h_3$\\
\midrule
cartesiano & $(x,y,z)$ & $(1,\,1,\,1)$ & $1$\\
cilíndrico & $(\rho,\varphi,z)$ & $(1,\,\rho,\,1)$ & $\rho$\\
esférico & $(r,\theta,\varphi)$ & $(1,\,r,\,r\sen\theta)$ & $r^{2}\sen\theta$\\
\bottomrule
\end{tabular}
\end{center}

% -------------------------------------------------------------------------
\subsection{Coordenadas cartesianas}\label{sec:cartesianas}
% -------------------------------------------------------------------------
\begin{figure}[ht]
\centering
\tdplotsetmaincoords{72}{120}
\begin{tikzpicture}[tdplot_main_coords, scale=1.45]
  \def\px{1.8}\def\py{2.3}\def\pz{2.1}
  \draw[eje] (0,0,0) -- (3.3,0,0) node[below left] {\small $x$};
  \draw[eje] (0,0,0) -- (0,3.6,0) node[right] {\small $y$};
  \draw[eje] (0,0,0) -- (0,0,3.2) node[above] {\small $z$};
  \fill[ocreteo, opacity=0.18] (0,0,0) -- (3.0,0,0) -- (3.0,3.2,0) -- (0,3.2,0) -- cycle;
  \draw[punteado] (\px,0,0) -- (\px,\py,0) -- (0,\py,0) (\px,\py,0) -- (\px,\py,\pz);
  \draw[punteado] (0,0,\pz) -- (\px,0,\pz) -- (\px,\py,\pz) (0,\py,0) -- (0,\py,\pz) -- (\px,\py,\pz);
  \draw[->, line width=1.5pt, azulprof] (0,0,0) -- (\px,\py,\pz);
  \node[azulprof, font=\small] at (0.55,0.5,1.6) {$\vc{r}$};
  \uvect{\px}{\py}{\pz}{(0.95,0,0)}{right}
  \uvect{\px}{\py}{\pz}{(0,0.95,0)}{right}
  \uvect{\px}{\py}{\pz}{(0,0,0.95)}{above}
  \node[rojoteo, font=\scriptsize] at (3.15,2.1,1.7) {$\ux$};
  \node[rojoteo, font=\scriptsize] at (2.0,3.15,1.7) {$\uy$};
  \node[rojoteo, font=\scriptsize] at (2.0,2.1,2.95) {$\uz$};
  \fill[azulprof] (\px,\py,\pz) circle (1.1pt);
  \node[font=\scriptsize] at (\px,-0.25,0) {$x$};
  \node[font=\scriptsize] at (-0.25,\py,0) {$y$};
  \node[font=\scriptsize] at (0,-0.15,\pz) {$z$};
\end{tikzpicture}
\caption{Coordenadas cartesianas $(x,y,z)$ y su base ortonormal constante $(\ux,\uy,\uz)$ en un punto $P$ de vector de posición $\vc{r}=x\,\ux+y\,\uy+z\,\uz$. Las líneas discontinuas muestran las proyecciones de $P$ sobre los ejes y los planos coordenados.}
\label{fig:coord-cart}
\end{figure}

Son las coordenadas $(x,y,z)$ asociadas a una base ortonormal \emph{constante}
$(\ux,\uy,\uz)$. Con $h_i=1$:
\begin{align}
  \vc{r}&=x\,\ux+y\,\uy+z\,\uz, & \dl&=\dif x\,\ux+\dif y\,\uy+\dif z\,\uz,\nonumber\\
  \dif S_x&=\dif y\,\dif z,\quad\dif S_y=\dif x\,\dif z,\quad\dif S_z=\dif x\,\dif y, & \dV&=\dif x\,\dif y\,\dif z.\nonumber
\end{align}
Los operadores son los ya conocidos:
\begin{align}
  \nab f&=\pd{f}{x}\,\ux+\pd{f}{y}\,\uy+\pd{f}{z}\,\uz, &
  \nab\cdot\vc{F}&=\pd{F_x}{x}+\pd{F_y}{y}+\pd{F_z}{z},\nonumber\\
  \nab\times\vc{F}&=\Bigl(\pd{F_z}{y}-\pd{F_y}{z}\Bigr)\ux+\Bigl(\pd{F_x}{z}-\pd{F_z}{x}\Bigr)\uy+\Bigl(\pd{F_y}{x}-\pd{F_x}{y}\Bigr)\uz, &
  \Lap f&=\pdd{f}{x}+\pdd{f}{y}+\pdd{f}{z}.\nonumber
\end{align}

% -------------------------------------------------------------------------
\subsection{Coordenadas cilíndricas}\label{sec:cilindricas}
% -------------------------------------------------------------------------
Se usan $(\rho,\varphi,z)$ con $\rho\ge0$ la distancia al eje $z$,
$\varphi\in[0,2\pi)$ el ángulo azimutal y $z\in\R$. (En el formulario del curso
la distancia al eje se denota $r$; aquí empleamos $\rho$ para no confundirla con
la distancia al origen $r$ de las coordenadas esféricas.)

\begin{figure}[ht]
\centering
\tdplotsetmaincoords{72}{120}
\begin{tikzpicture}[tdplot_main_coords, scale=1.45]
  \def\rh{2.3}\def\ph{38}\def\pz{1.7}
  \pgfmathsetmacro{\px}{\rh*cos(\ph)}\pgfmathsetmacro{\py}{\rh*sin(\ph)}
  \draw[eje] (0,0,0) -- (3.3,0,0) node[below left] {\small $x$};
  \draw[eje] (0,0,0) -- (0,3.6,0) node[right] {\small $y$};
  \draw[eje] (0,0,0) -- (0,0,3.2) node[above] {\small $z$};
  % pared cilíndrica
  \fill[ocreteo, opacity=0.20]
    plot[variable=\t, domain=-15:105, samples=14] ({\rh*cos(\t)},{\rh*sin(\t)},-0.3)
    -- plot[variable=\t, domain=105:-15, samples=14] ({\rh*cos(\t)},{\rh*sin(\t)},2.6) -- cycle;
  \draw[ocreteo!70, line width=0.4pt] plot[variable=\t, domain=-15:105, samples=14] ({\rh*cos(\t)},{\rh*sin(\t)},\pz);
  \draw[punteado] (0,0,0) -- (\px,\py,0) -- (\px,\py,\pz);
  \draw[->, line width=1.5pt, azulprof] (0,0,0) -- (\px,\py,\pz);
  \node[azulprof, font=\small] at (0.35,0.1,1.45) {$\vc{r}$};
  \draw[grispizarra, line width=0.5pt] plot[variable=\t, domain=0:\ph, samples=8] ({0.9*cos(\t)},{0.9*sin(\t)},0);
  \node[font=\scriptsize] at ({1.25*cos(\ph/2)},{1.25*sin(\ph/2)},0) {$\varphi$};
  \node[font=\scriptsize] at ({0.72*\rh*cos(\ph)+0.30*sin(\ph)},{0.72*\rh*sin(\ph)-0.30*cos(\ph)},0) {$\rho$};
  \node[font=\scriptsize] at (\px+0.05,\py+0.25,0.5*\pz) {$z$};
  \uvect{\px}{\py}{\pz}{({0.95*cos(\ph)},{0.95*sin(\ph)},0)}{right}
  \uvect{\px}{\py}{\pz}{({-0.95*sin(\ph)},{0.95*cos(\ph)},0)}{above}
  \uvect{\px}{\py}{\pz}{(0,0,0.95)}{above}
  \node[rojoteo, font=\scriptsize] at ({\px+1.25*cos(\ph)},{\py+1.25*sin(\ph)},\pz-0.1) {$\urho$};
  \node[rojoteo, font=\scriptsize] at ({\px-1.2*sin(\ph)},{\py+1.2*cos(\ph)},\pz+0.05) {$\uph$};
  \node[rojoteo, font=\scriptsize] at (\px,\py,\pz+1.25) {$\uz$};
  \fill[azulprof] (\px,\py,\pz) circle (1.8pt);
\end{tikzpicture}
\caption{Coordenadas cilíndricas $(\rho,\varphi,z)$. El versor $\urho$ apunta alejándose del eje $z$ y $\uph$ es tangente a la circunferencia de radio $\rho$ (en el sentido en que crece $\varphi$); $\uz$ es el mismo que en cartesianas. La pared cilíndrica es la superficie $\rho=\text{cte}$. Al cambiar de punto, $\urho$ y $\uph$ giran.}
\label{fig:coord-cil}
\end{figure}


\begin{align}
  x&=\rho\cos\varphi, & y&=\rho\sen\varphi, & z&=z,\nonumber\\
  \rho&=\sqrt{x^{2}+y^{2}}, & \varphi&=\tan^{-1}(y/x), & z&=z .\nonumber
\end{align}
Los versores y su relación con la base cartesiana son
\begin{align}
  \urho&=\cos\varphi\,\ux+\sen\varphi\,\uy, & \ux&=\cos\varphi\,\urho-\sen\varphi\,\uph,\nonumber\\
  \uph&=-\sen\varphi\,\ux+\cos\varphi\,\uy, & \uy&=\sen\varphi\,\urho+\cos\varphi\,\uph,\label{eq:base-cil}\\
  \uz&=\uz, & \uz&=\uz.\nonumber
\end{align}
La posición es $\vc{r}=\rho\,\urho+z\,\uz$ y los elementos geométricos
(\cref{prop:elementos} con $h=(1,\rho,1)$):
\begin{align}
  \dl&=\dif\rho\,\urho+\rho\,\dif\varphi\,\uph+\dif z\,\uz, \nonumber\\
  \dif S_\rho&=\rho\,\dif\varphi\,\dif z,\quad
  \dif S_\varphi=\dif\rho\,\dif z,\quad
  \dif S_z=\rho\,\dif\rho\,\dif\varphi,\qquad
  \dV=\rho\,\dif\rho\,\dif\varphi\,\dif z .\nonumber
\end{align}
Los operadores (\cref{teo:operadores-ortogonales}) son
\begin{align}
  \nab f&=\pd{f}{\rho}\,\urho+\frac{1}{\rho}\pd{f}{\varphi}\,\uph+\pd{f}{z}\,\uz,\nonumber\\
  \nab\cdot\vc{F}&=\frac{1}{\rho}\pd{(\rho F_\rho)}{\rho}+\frac{1}{\rho}\pd{F_\varphi}{\varphi}+\pd{F_z}{z},\nonumber\\
  \nab\times\vc{F}&=\Bigl(\frac{1}{\rho}\pd{F_z}{\varphi}-\pd{F_\varphi}{z}\Bigr)\urho
    +\Bigl(\pd{F_\rho}{z}-\pd{F_z}{\rho}\Bigr)\uph
    +\frac{1}{\rho}\Bigl(\pd{(\rho F_\varphi)}{\rho}-\pd{F_\rho}{\varphi}\Bigr)\uz,\label{eq:ops-cil}\\
  \Lap f&=\frac{1}{\rho}\pd{}{\rho}\Bigl(\rho\pd{f}{\rho}\Bigr)+\frac{1}{\rho^{2}}\pdd{f}{\varphi}+\pdd{f}{z}.\nonumber
\end{align}

% -------------------------------------------------------------------------
\subsection{Coordenadas esféricas}\label{sec:esfericas}
% -------------------------------------------------------------------------
Se usan $(r,\theta,\varphi)$ con $r\ge0$ la distancia al origen,
$\theta\in[0,\pi]$ el ángulo polar (medido desde el eje $z$) y
$\varphi\in[0,2\pi)$ el azimut. Obsérvese que $\sen\theta\ge0$ en todo el
rango, hecho que usaremos al calcular jacobianos.

\begin{figure}[ht]
\centering
\tdplotsetmaincoords{72}{120}
\begin{tikzpicture}[tdplot_main_coords, scale=1.45]
  \def\rr{3.0}\def\th{52}\def\ph{38}
  \pgfmathsetmacro{\px}{\rr*sin(\th)*cos(\ph)}\pgfmathsetmacro{\py}{\rr*sin(\th)*sin(\ph)}\pgfmathsetmacro{\pz}{\rr*cos(\th)}
  \draw[eje] (0,0,0) -- (3.5,0,0) node[below left] {\small $x$};
  \draw[eje] (0,0,0) -- (0,3.7,0) node[right] {\small $y$};
  \draw[eje] (0,0,0) -- (0,0,3.5) node[above] {\small $z$};
  % casquete esférico
  \fill[ocreteo, opacity=0.20]
    plot[variable=\t, domain=0:90, samples=10] ({\rr*sin(\t)*cos(-10)},{\rr*sin(\t)*sin(-10)},{\rr*cos(\t)})
    -- plot[variable=\t, domain=-10:100, samples=14] ({\rr*cos(\t)},{\rr*sin(\t)},0)
    -- plot[variable=\t, domain=90:0, samples=10] ({\rr*sin(\t)*cos(100)},{\rr*sin(\t)*sin(100)},{\rr*cos(\t)}) -- cycle;
  \draw[punteado] (0,0,0) -- (\px,\py,0) -- (\px,\py,\pz);
  \draw[->, line width=1.5pt, azulprof] (0,0,0) -- (\px,\py,\pz);
  \node[azulprof, font=\small] at ({0.45*\px+0.25},{0.45*\py-0.1},{0.45*\pz+0.25}) {$\vc{r}$};
  % arcos de theta y phi
  \draw[grispizarra, line width=0.5pt] plot[variable=\t, domain=0:\th, samples=8]
        ({1.0*sin(\t)*cos(\ph)},{1.0*sin(\t)*sin(\ph)},{1.0*cos(\t)});
  \node[font=\scriptsize] at ({1.35*sin(\th/2)*cos(\ph)},{1.35*sin(\th/2)*sin(\ph)},{1.35*cos(\th/2)}) {$\theta$};
  \draw[grispizarra, line width=0.5pt] plot[variable=\t, domain=0:\ph, samples=8] ({1.0*cos(\t)},{1.0*sin(\t)},0);
  \node[font=\scriptsize] at ({1.35*cos(\ph/2)},{1.35*sin(\ph/2)},0) {$\varphi$};
  \uvect{\px}{\py}{\pz}{({0.95*sin(\th)*cos(\ph)},{0.95*sin(\th)*sin(\ph)},{0.95*cos(\th)})}{above}
  \uvect{\px}{\py}{\pz}{({0.95*cos(\th)*cos(\ph)},{0.95*cos(\th)*sin(\ph)},{-0.95*sin(\th)})}{below}
  \uvect{\px}{\py}{\pz}{({-0.95*sin(\ph)},{0.95*cos(\ph)},0)}{right}
  \node[rojoteo, font=\scriptsize] at ({\px+1.2*sin(\th)*cos(\ph)},{\py+1.2*sin(\th)*sin(\ph)},{\pz+1.2*cos(\th)}) {$\ur$};
  \node[rojoteo, font=\scriptsize] at ({\px+1.25*cos(\th)*cos(\ph)},{\py+1.25*cos(\th)*sin(\ph)},{\pz-1.25*sin(\th)}) {$\uth$};
  \node[rojoteo, font=\scriptsize] at ({\px-1.2*sin(\ph)},{\py+1.2*cos(\ph)},\pz) {$\uph$};
  \fill[azulprof] (\px,\py,\pz) circle (1.1pt);
\end{tikzpicture}
\caption{Coordenadas esféricas $(r,\theta,\varphi)$. El versor $\ur$ es radial, $\uth$ apunta hacia la dirección en que aumenta $\theta$ (hacia el sur) y $\uph$ es el mismo versor azimutal de las cilíndricas. El casquete es la superficie $r=\text{cte}$; el vector de posición es $\vc{r}=r\,\ur$.}
\label{fig:coord-esf}
\end{figure}

\begin{align}
  x&=r\sen\theta\cos\varphi, & y&=r\sen\theta\sen\varphi, & z&=r\cos\theta,\nonumber\\
  r&=\sqrt{x^{2}+y^{2}+z^{2}}, & \theta&=\tan^{-1}\!\Bigl(\sqrt{x^{2}+y^{2}}/z\Bigr), & \varphi&=\tan^{-1}(y/x).\nonumber
\end{align}
Los versores esféricos en función de los cartesianos, y viceversa, son
\begin{align}
  \ur&=\sen\theta\cos\varphi\,\ux+\sen\theta\sen\varphi\,\uy+\cos\theta\,\uz,\nonumber\\
  \uth&=\cos\theta\cos\varphi\,\ux+\cos\theta\sen\varphi\,\uy-\sen\theta\,\uz,\nonumber\\
  \uph&=-\sen\varphi\,\ux+\cos\varphi\,\uy,\label{eq:base-esf}\\[4pt]
  \ux&=\sen\theta\cos\varphi\,\ur+\cos\theta\cos\varphi\,\uth-\sen\varphi\,\uph,\nonumber\\
  \uy&=\sen\theta\sen\varphi\,\ur+\cos\theta\sen\varphi\,\uth+\cos\varphi\,\uph,\nonumber\\
  \uz&=\cos\theta\,\ur-\sen\theta\,\uth.\nonumber
\end{align}
La posición es simplemente $\vc{r}=r\,\ur$ (¡la dirección está en $\ur$!) y
\begin{align}
  \dl&=\dif r\,\ur+r\,\dif\theta\,\uth+r\sen\theta\,\dif\varphi\,\uph, \nonumber\\
  \dif S_r&=r^{2}\sen\theta\,\dif\theta\,\dif\varphi,\quad
  \dif S_\theta=r\sen\theta\,\dif r\,\dif\varphi,\quad
  \dif S_\varphi=r\,\dif r\,\dif\theta,\qquad
  \dV=r^{2}\sen\theta\,\dif r\,\dif\theta\,\dif\varphi .\nonumber
\end{align}
Los operadores son
\begin{align}
  \nab f&=\pd{f}{r}\,\ur+\frac{1}{r}\pd{f}{\theta}\,\uth+\frac{1}{r\sen\theta}\pd{f}{\varphi}\,\uph,\nonumber\\
  \nab\cdot\vc{F}&=\frac{1}{r^{2}}\pd{(r^{2}F_r)}{r}+\frac{1}{r\sen\theta}\pd{(F_\theta\sen\theta)}{\theta}+\frac{1}{r\sen\theta}\pd{F_\varphi}{\varphi},\nonumber\\
  \nab\times\vc{F}&=\frac{1}{r\sen\theta}\Bigl(\pd{(F_\varphi\sen\theta)}{\theta}-\pd{F_\theta}{\varphi}\Bigr)\ur
    +\frac{1}{r}\Bigl(\frac{1}{\sen\theta}\pd{F_r}{\varphi}-\pd{(rF_\varphi)}{r}\Bigr)\uth
    +\frac{1}{r}\Bigl(\pd{(rF_\theta)}{r}-\pd{F_r}{\theta}\Bigr)\uph,\label{eq:ops-esf}\\
  \Lap f&=\frac{1}{r^{2}}\pd{}{r}\Bigl(r^{2}\pd{f}{r}\Bigr)+\frac{1}{r^{2}\sen\theta}\pd{}{\theta}\Bigl(\sen\theta\pd{f}{\theta}\Bigr)+\frac{1}{r^{2}\sen^{2}\theta}\pdd{f}{\varphi}.\nonumber
\end{align}

\begin{ejemplo}[Campos con simetría esférica]\label{ej:simetria-esferica}
Si $\vc{A}=A(r)\,\ur$ solo tiene componente radial y depende solo de $r$, las
\cref{eq:ops-esf} dan
\[
  \nab\cdot\vc{A}=\frac{1}{r^{2}}\frac{\dif}{\dif r}\bigl(r^{2}A\bigr),
  \qquad\nab\times\vc{A}=\vc{0},
\]
y para un campo escalar $f(r)$, $\nab f=f'(r)\,\ur$ y
$\Lap f=\frac{1}{r^{2}}\frac{\dif}{\dif r}\bigl(r^{2}f'\bigr)$. En particular,
$A=r^{n}$ da $\nab\cdot\vc{A}=(n+2)\,r^{n-1}$ y, para $A=1/r^{2}$ (campo de una
carga puntual), la divergencia se anula \emph{fuera del origen}.
\end{ejemplo}

% -------------------------------------------------------------------------
\subsection{Derivadas de los versores y jacobianos}\label{sec:versores-jacobianos}
% -------------------------------------------------------------------------
En cilíndricas y esféricas los versores \emph{dependen del punto}, y derivar un
campo exige derivar también la base. Se obtienen directamente de
las \cref{eq:base-cil,eq:base-esf}.

\begin{lema}[Derivadas de los versores]\label{lema:derivadas-versores}
Las únicas derivadas no nulas son, en cilíndricas,
\[
  \pd{\urho}{\varphi}=\uph,\qquad\pd{\uph}{\varphi}=-\urho,
\]
y en esféricas,
\[
  \pd{\ur}{\theta}=\uth,\quad\pd{\ur}{\varphi}=\sen\theta\,\uph,\quad
  \pd{\uth}{\theta}=-\ur,\quad\pd{\uth}{\varphi}=\cos\theta\,\uph,\quad
  \pd{\uph}{\varphi}=-\sen\theta\,\ur-\cos\theta\,\uth .
\]
\end{lema}
\begin{proof}
Se derivan los desarrollos cartesianos \eqref{eq:base-cil} y \eqref{eq:base-esf} (con
$\ux,\uy,\uz$ constantes) y se reexpresa el resultado en la base local. Por
ejemplo, $\partial_\varphi\urho=-\sen\varphi\,\ux+\cos\varphi\,\uy=\uph$ y
$\partial_\theta\ur=\cos\theta\cos\varphi\,\ux+\cos\theta\sen\varphi\,\uy-\sen\theta\,\uz=\uth$.
\end{proof}

Estas relaciones explican, por ejemplo, por qué el laplaciano vectorial no es
el laplaciano de cada componente cuando se usan curvilíneas
(\cref{def:laplaciano}), y por qué aparecen los factores $\rho$ y $r\sen\theta$
en los operadores.

\begin{proposicion}[Jacobianos]\label{prop:jacobianos}
Para el cambio de variable a coordenadas cilíndricas y esféricas,
\[
  \Bigl|\frac{\partial(x,y,z)}{\partial(\rho,\varphi,z)}\Bigr|=\rho,
  \qquad
  \Bigl|\frac{\partial(x,y,z)}{\partial(r,\theta,\varphi)}\Bigr|=r^{2}\sen\theta,
\]
de modo que $\iiint_Vf\,\dif x\,\dif y\,\dif z$ se transforma, respectivamente, en
\[
  \iiint f\,\rho\,\dif\rho\,\dif\varphi\,\dif z
  \qquad\text{y en}\qquad
  \iiint f\,r^{2}\sen\theta\,\dif r\,\dif\theta\,\dif\varphi .
\]
\end{proposicion}
\begin{proof}
Es el caso particular $|\det|=h_1h_2h_3$ de la \cref{prop:elementos} (el
valor absoluto elimina el signo, y en esféricas $\sen\theta\ge0$ porque
$\theta\in[0,\pi]$). El determinante explícito en esféricas se calcula, paso a
paso, en el \cref{ejer:jacobiano}.
\end{proof}
